%!TEX root = ../main.tex

\chapter{Dataset and Evaluation Framework}\label{ch:dataset}

This chapter describes the dataset used in all experiments, the evaluation metrics, and
the fixed subset of 500 examples that was used for most comparisons during development.


% ===========================================================================
\section{MultiDoc2Dial}\label{sec:dataset:mdd}

\Gls{mdd}~\cite{feng2021multidoc2dial} is a dataset of goal-oriented dialogues in which
each agent turn is grounded in one or more spans of a collection of government-service
documents. It was created to study how systems retrieve and use evidence from several
documents during a conversation.

\subsection{Corpus}

The documents come from four US government-service websites: the New York Department of
Motor Vehicles (DMV: driver's licences, vehicle registration, traffic tickets), the
Social Security Administration (SSA: retirement, disability and survivor benefits),
Federal Student Aid (StudentAid: financial aid and student loans) and Veterans Affairs
(VA: health care and benefits for veterans).

The corpus has 488 documents with 35{,}659 annotated spans, which were grouped into
4{,}114 sections. A section is a heading together with the text that follows it, and it
is the retrieval unit in all experiments.

\subsection{Dialogues and Annotations}

The dialogues alternate between a user and an agent. They were collected through
crowd-sourcing: annotators asked the kind of questions a citizen would ask a
government-service agent, and the agent answers were grounded in specific document spans.
Table~\ref{tab:dataset:split} gives the size of the three splits.

\begin{table}[ht]
    \centering
    \caption{MultiDoc2Dial splits.}
    \label{tab:dataset:split}
    \begin{tabular}{lrrrr}
        \toprule
        Split & Dialogues & Questions & Unanswerable & Doc switch \\
        \midrule
        Train      & 3{,}474 & 22{,}806 & 1{,}355 (5.9\%) & 4{,}392 (19.3\%) \\
        Validation &   661   &  4{,}427 &   226 (5.1\%)   &   905 (20.4\%) \\
        Test       &   661   &  4{,}335 &   241 (5.6\%)   &   736 (17.0\%) \\
        \bottomrule
    \end{tabular}
\end{table}

Each agent turn is annotated with its \emph{grounding spans}, the parts of the documents
the answer comes from. These annotations are used to measure retrieval (Recall@$K$) and to
identify unanswerable turns. No model is trained in this work; the training split is only
used to select the few-shot examples of the generation prompt.

\subsection{Unanswerable and Document-Switch Turns}

Two kinds of turns are evaluated separately. About 5\% of the questions are
\emph{unanswerable}: the documents do not contain the answer, and the reference answer is
a refusal. A turn is a \emph{document switch} when its gold document differs from the one
of the previous agent turn in the same dialogue; this is the case for about a fifth of
the validation questions. These turns are harder for retrieval, because the recent
dialogue is about a different document than the question. They are analysed separately
in Section~\ref{sec:rag:fullsplit}.


% ===========================================================================
\section{Evaluation Metrics}\label{sec:dataset:metrics}

\subsection{Token F1}

Answers are evaluated with \emph{token \fone{}} against the reference answer, after
lowercasing and removing punctuation and articles, as in SQuAD~\cite{rajpurkar2016squad}:
\begin{equation}
    \fone = 2 \cdot \frac{\text{precision} \cdot \text{recall}}
                         {\text{precision} + \text{recall}},
\end{equation}
where precision and recall are computed over the words the answer and the reference have
in common. Token \fone{} gives partial credit for partly matching answers, unlike exact
match (\gls{em}), but it still only measures word overlap: a correct answer worded
differently from the reference gets a low score.

\subsection{Recall@K}

Retrieval is measured by \emph{Recall@$K$}, the fraction of gold grounding spans that are
contained in the top-$K$ retrieved sections. For a question with grounding set $G$,
\begin{equation}
    \ratk{K} = \frac{|\{g \in G : g \in \text{top-}K\}|}{|G|}.
\end{equation}
$\ratk{1}$ and $\ratk{5}$ are reported. Since five sections are given to the generator in
all experiments, $\ratk{5}$ says whether the model had access to the relevant text.

\subsection{Abstention Metrics}

The model is instructed to answer ``Unfortunately, no relevant information is found.''
when the context does not contain the answer, and any answer containing ``no relevant
information'' counts as an abstention (see Appendix~\ref{ax:log:pilot} for an earlier,
stricter criterion). Three metrics describe how well this works. Abstain recall is the
fraction of unanswerable questions on which the model abstains, the false-abstention rate
is the fraction of answerable questions on which it abstains, and abstain precision is the
fraction of abstentions that fall on unanswerable questions. A good system has high
abstain recall and a low false-abstention rate; the two pull in opposite directions, since
a model that always abstains has perfect abstain recall but answers nothing.


% ===========================================================================
\section{Dev-500 Subset}\label{sec:dataset:dev500}

Running one pipeline on the full validation split (4{,}427 questions) takes several hours
of \gls{llm} calls, which is too slow for the many comparisons needed during development.
A fixed subset of 500 validation questions, called \emph{dev-500}, was therefore used for
most experiments.

The subset is balanced across domains: 125 questions were drawn at random from each
domain with \texttt{random.Random(seed=42)}. At least 30 unanswerable questions were
required, and the random draw already contained 31, so nothing had to be added. The ids
are stored in \texttt{data/dev\_500\_ids.json}, so every experiment uses exactly the same
questions. Table~\ref{tab:dataset:dev500} shows its composition.

\begin{table}[ht]
    \centering
    \caption{Composition of the dev-500 subset.}
    \label{tab:dataset:dev500}
    \begin{tabular}{lrrr}
        \toprule
        Domain & Questions & Unanswerable & Doc switch \\
        \midrule
        DMV          & 125 &  4 & 21 \\
        SSA          & 125 &  9 & 18 \\
        StudentAid   & 125 & 13 & 31 \\
        VA           & 125 &  5 & 24 \\
        \midrule
        Total        & 500 & 31 (6.2\%) & 94 (18.8\%) \\
        \bottomrule
    \end{tabular}
\end{table}

Dev-500 is used for the prompt and reranker experiments
(Sections~\ref{sec:rag:prompt} and~\ref{sec:rag:reranker}), the oracle analysis
(Section~\ref{sec:rag:oracle}) and all knowledge-graph retrieval experiments. The five
base pipelines are also reported on the full validation split
(Section~\ref{sec:rag:fullsplit}), and the final systems on the full validation and the
test split (Chapter~\ref{ch:final}).
